\documentclass{article}
\usepackage{amsmath,amssymb,booktabs}
\usepackage[margin=1in]{geometry}
\usepackage[colorlinks=true,linkcolor=blue,citecolor=blue,urlcolor=blue]{hyperref}

\newcommand{\tr}{\operatorname{Tr}}
\newcommand{\Da}{\Delta a}
\newcommand{\mpi}{m_\pi}

\title{$A'A'\to\gamma\gamma$ and $A'A'\to M\bar M$ for the axial-vector mediator}
\author{\texttt{AA\_sm\_xsec.m}, \texttt{sv\_AAsm\_grid.m}, \texttt{freeze\_coupled.m} ($c_{AA}$)}
\date{2026-11-17}

\begin{document}
\maketitle

\section{Why}

For $m_{A'}<\min(2m_\chi,\,3\mpi)$ the $A'$ cannot decay. Landau--Yang forbids $A'\to\gamma\gamma$,
$A'\to\pi^0\gamma$ and $A'\to3\gamma$ are $C$-odd (the axial current is $C$-even), and $A'$ has no
lepton couplings. Its only number-changing link was $\chi\bar\chi\leftrightarrow A'A'$, so
\texttt{freeze\_coupled} left an $A'$ relic that was never counted. Here we add the lowest-order
$A'A'\to$ SM channels:
\begin{itemize}
\item $A'A'\to\gamma\gamma$ through one charged-meson loop;
\item $A'A'\to\pi^+\pi^-,\,K^+K^-,\,K^0\bar K^0$ at tree level, which is the same vertex.
\end{itemize}
We also return $\Omega_{A'}h^2$ so that overclosure can be flagged.

\section{The vertex}

The axial source is $a_\mu=A'_\mu\,\mathrm{diag}(a_u,a_d,a_s)$, with
$U=\exp(i\sqrt2\Phi/F)$ and $D_\mu U=\partial_\mu U-i\{a_\mu,U\}$. The term of
$\tfrac{F^2}{4}\tr(D_\mu UD^\mu U^\dagger)$ quadratic in $a$ is
$\tfrac{F^2}{4}\tr(\{a,U\}\{a,U^\dagger\})$. With $X=\sqrt2\Phi/F$, its $O(X^2)$ part is
$\tfrac{F^2}{4}\big[2\tr(aXaX)-2\tr(a^2X^2)\big]=\tr(a\Phi a\Phi-a^2\Phi^2)$. Because $a$ is
diagonal,
\begin{equation}
\tr(a\Phi a\Phi-a^2\Phi^2)=\sum_{i,j}(a_ia_j-a_i^2)\Phi_{ij}\Phi_{ji}
=-\sum_{i<j}(a_i-a_j)^2\,\Phi_{ij}\Phi_{ji}.
\end{equation}
This gives
\begin{equation}
\boxed{\mathcal L_{aa}=-A'\!\cdot\!A'\Big[\Da_\pi^2\,\pi^+\pi^-+\Da_K^2\,K^+K^-+\Da_{K^0}^2\,K^0\bar K^0\Big]},
\qquad \Da_\pi=a_u-a_d,\ \Da_K=a_u-a_s,\ \Da_{K^0}=a_d-a_s .
\end{equation}
This is the contact of \texttt{notes/threepi\_annihilation.tex} (section on $\pi^+\pi^-A'$) with the kaons
added. The diagonal mesons ($\pi^0,\eta,\eta'$) commute with the source and drop out. The
Feynman rule is $-2i\Da^2(\varepsilon_1\!\cdot\!\varepsilon_2)$, where the 2 comes from the two
identical $A'$.

\paragraph{Which diagrams survive at one loop.} The candidates are:
\begin{enumerate}
\item The contact $\times$ a scalar-QED loop (two $M\bar M\gamma$ vertices, or one
  $M\bar M\gamma\gamma$ seagull). This is the one that contributes.
\item $A'\to\pi^0$ mixing on an external leg. This is $\propto k\cdot\varepsilon=0$ on shell, so it
  vanishes.
\item $A'\gamma M$ vertices. There are none: $C$ and $P$ forbid them, and $\tr(v\ldots a\ldots)$
  cross terms vanish for commuting diagonal sources.
\item An $O(p^4)$ $A'A'\gamma\gamma$ counterterm. $L_{10}\tr(U^\dagger F_R UF_L)$ has none at this
  order in $\Phi$, so the loop is finite.
\end{enumerate}
The $A'A'$ pair therefore enters exactly like a scalar source $\sigma=A'\!\cdot\!A'$ coupled to
$|M|^2$, i.e.\ like $h\to\gamma\gamma$ through a charged scalar, with $m_h^2\to s$.

\section{The loop}

The effective Lagrangian is
\begin{equation}
\mathcal L_{\rm eff}=c(s)\,(A'\!\cdot\!A')\,F_{\mu\nu}F^{\mu\nu},\qquad
c(s)=\frac{\alpha}{16\pi}\sum_{j=\pi^\pm,K^\pm}\frac{\Da_j^2}{m_j^2}\,A_0\!\Big(\frac{s}{4m_j^2}\Big),
\end{equation}
with
\begin{equation}
A_0(\tau)=\frac{f(\tau)-\tau}{\tau^2},\quad
f(\tau)=\begin{cases}\arcsin^2\sqrt\tau & \tau\le1\\[2pt]
-\tfrac14\Big[\ln\frac{1+\sqrt{1-1/\tau}}{1-\sqrt{1-1/\tau}}-i\pi\Big]^2 & \tau>1\end{cases},
\qquad A_0=\tfrac13+\tfrac{8\tau}{45}+\tfrac{4\tau^2}{35}+\dots
\end{equation}
The overall sign is a convention. Both mesons enter with the same sign, because $\Da_j^2>0$ and
$A_0>0$ below threshold, so the $\pi$ and $K$ loops add.

\paragraph{Check 1: low-energy theorem.} At $s\to0$ the contact is a mass shift
$\delta m_j^2=\Da_j^2A'\!\cdot\!A'$. A charged scalar adds
$-\tfrac14\tfrac{e^2}{48\pi^2}\ln m^2\,F^2$ to the photon kinetic term. Varying it gives
$\mp\tfrac{\alpha}{48\pi}\tfrac{\delta m^2}{m^2}F^2=\mp\tfrac{\alpha}{16\pi}\tfrac{\Da^2}{m^2}\tfrac13\,A'\!\cdot\!A'F^2$,
which is $c(0)$ with $A_0(0)=\tfrac13$.

\paragraph{Check 2: $h\to\gamma\gamma$.} Taking $\tfrac14 g\,hFF$ in place of $c\,A'A'FF$ gives
$\Gamma=g^2m^3/(64\pi)$, the standard normalisation.

\section{Cross sections}

\paragraph{Amplitude.} $F_{\mu\nu}F^{\mu\nu}$ between two photons gives
$4[(k_1\!\cdot\!k_2)(e_1\!\cdot\!e_2)-(k_1\!\cdot\!e_2)(k_2\!\cdot\!e_1)]$, and
$A'\!\cdot\!A'$ gives $2\,\varepsilon_1\!\cdot\!\varepsilon_2$. So
$\mathcal M=8c\,(\varepsilon_1\!\cdot\!\varepsilon_2)[\dots]$. Summing over photon
polarisations, $\sum|[\dots]|^2=2(k_1\!\cdot\!k_2)^2=s^2/2$.

\paragraph{A$'$ polarisations.} With $-g+kk/m_A^2$ and $p_1\!\cdot\!p_2=(s-2m_A^2)/2$:
\begin{equation}
P(s)\equiv\sum_{\rm pol}|\varepsilon_1\!\cdot\!\varepsilon_2|^2=2+\frac{(s-2m_A^2)^2}{4m_A^4}.
\end{equation}
The second term is the two longitudinal $A'$. It is physical: by Goldstone equivalence
$\varepsilon_L\simeq p/m_A$ is the $\pi^0$ admixture, with angle $\sim F\Da/m_A$. It makes both
channels grow as $s^2/m_A^4$ for a light $A'$. This is the same $1/m_A^2$ as in
$\pi^+\pi^-A'$, and it fails where the small-mixing expansion does.

\paragraph{Results.} The initial-state average is $\tfrac19$, there is a factor $\tfrac12$ for
identical photons, and $\beta_A=\sqrt{1-4m_A^2/s}$. Then
\begin{align}
\sum|\mathcal M_{\gamma\gamma}|^2&=32|c|^2s^2P(s), &
\sigma_{\gamma\gamma}&=\frac{|c(s)|^2\,s\,P(s)}{9\pi\,\beta_A},\\
\sum|\mathcal M_{M\bar M}|^2&=4\Da_j^4P(s), &
\sigma_{M\bar M}&=\frac{\Da_j^4\,P(s)\,\beta_M}{36\pi\,s\,\beta_A}.
\end{align}
\paragraph{Non-relativistic limit.} Take $v_{\rm rel}=2\beta_A$, $s\to4m_A^2$ and $P\to3$. With
the pion only and $A_0=\tfrac13$,
\begin{equation}
\sigma_{\gamma\gamma}v_{\rm rel}\to\frac{8|c|^2m_A^2}{3\pi}
=\frac{\alpha^2\Da_\pi^4m_A^2}{864\pi^3\mpi^4}.
\end{equation}
This is $s$-wave, so it survives late. The code reproduces it, and the ratio to the LO pion-only
value is $[A_0(m_A^2/\mpi^2)/A_0(0)\,(1+\text{kaon})]^2=(1.419\times1.058)^2=2.25$ at
$m_A=0.1$, $a\propto[2,-1,-1]$.

\paragraph{Order in $\chi$PT.} This is LO $\Da$. A one-loop graph carries the $L_4$ factor
\texttt{corr4} to the power $V-I=0$, so it is not rescaled. The tree-level $3\pi$ and
$\pi\pi A'$ currents are rescaled (\texttt{threepi\_annihilation.tex}). The $O(p^4)$ corrections
to the contact are not included; expect tens of percent.

\section{Boltzmann term and relic}

$A'$ is self-conjugate. Each annihilation removes two particles, which cancels the $\tfrac12$
for identical initial states, so
\begin{equation}
\frac{dY_{A'}}{d\ln x}\supset-\frac{s\langle\sigma v\rangle_{AA}}{H}\,(Y_{A'}^2-Y_{A',\rm eq}^2),
\qquad c_{AA}=2\lambda\,\langle\sigma v\rangle_{AA}.
\end{equation}
Here $\lambda$ is \texttt{freeze\_coupled}'s coefficient. It carries the Dirac $\tfrac12$ for
$\chi\bar\chi$, hence the 2. $\langle\sigma v\rangle$ is the Gondolo--Gelmini average with
Maxwell--Boltzmann statistics (matching $Y_{\rm eq}$), $g_{A'}=3$, tabulated on the solver's $T$
grid. It is piecewise between the thresholds $2m_A,2m_{\pi^\pm},2m_{K^\pm},2m_{K^0}$
(\texttt{sv\_AAsm\_grid.m}).

\paragraph{Where it enters.} It is added to every path in \texttt{freeze\_coupled}:
\begin{itemize}
\item the full $2\times2$ system and its Jacobians;
\item the dark-exchange shortcut, where it is one of $N$'s slow links;
\item the fast-$\phi$ quasi-steady state. There the quadratic coefficient becomes
  $c_{\rm rev}+c_{AA}$, $D'\mathrel{+}=2c_{AA}Y_{\rm eq}$, and the exchange term $T_\phi$ gains
  $-2c_{AA}(2Y_{\rm eq}\delta+\delta^2)$;
\item the $Y_{A'}$ initial condition, as a rate $c_{AA}Y_{\rm eq}$.
\end{itemize}
\texttt{LOOP\_AA\_SM=0} turns it off.

\paragraph{Second output.} \texttt{[Oh2, info] = freeze\_coupled(...)} returns:
\begin{itemize}
\item $Y_{A'}$ at $x_{\rm final}$;
\item $\Omega_{A'}h^2=m_A s_0Y_{A'}/(\rho_c/h^2)$;
\item $\Gamma_{A'}$;
\item \texttt{phi\_stable} ($\Gamma\,t_0<1$);
\item \texttt{overclosed} (stable and $\Omega_{A'}h^2>0.12$).
\end{itemize}
$\Omega_\chi h^2$ is unchanged in meaning. Whether to add $\Omega_{A'}$ to the relic target is
left open.

\section{$A'\pi\to\pi\pi$}

The term of the axial current linear in $a$, to $O(\Phi^3)$, is
(\texttt{threepi\_annihilation.tex})
\begin{equation}
\mathcal L_a=-\Da\,A'_\mu\Big[F\partial^\mu\pi^3+\frac{2}{3F}\big(\pi^3\,\pi\!\cdot\!\partial^\mu\pi
-\pi^2\partial^\mu\pi^3\big)\Big]
\supset-\frac{2\Da}{3F}A'_\mu\Big[\pi^0\partial^\mu(\pi^+\pi^-)-2\pi^+\pi^-\partial^\mu\pi^0\Big].
\end{equation}
In $A'\to\pi^+\pi^-\pi^0$ the current is
$V^\mu=F q^\mu M_4/(s-m_{\pi^0}^2)-\tfrac{2}{3F}(p_++p_--2p_0)^\mu$, with $q=p_++p_-+p_0$. Crossed
to $2\to2$, the pole piece is $\propto k\cdot\varepsilon=0$ for an on-shell $A'$. Only the
contact survives:
\begin{itemize}
\item $A'(k)\pi^0(p)\to\pi^+\pi^-$: $p_0\to-p$, so $V=\tfrac{2}{3F}(p_++p_-+2p)=\tfrac{2}{3F}(k+3p)$
  and $\mathcal M=\tfrac{2\Da}{F}\,\varepsilon\cdot p$.
\item $A'(k)\pi^\pm(p)\to\pi^\pm(p')\pi^0(p_0)$: $V\propto p'-p-2p_0=k-3p_0$, so
  $\mathcal M=-\tfrac{2\Da}{F}\,\varepsilon\cdot p_0$.
\end{itemize}
Summing over the $A'$ polarisations,
$\sum|\varepsilon\cdot p|^2=(k\cdot p)^2/m_A^2-p^2$. This gives
\begin{align}
\sum_{\rm pol}|\mathcal M_{0\to+-}|^2&=\frac{4\Da^2}{F^2}\,\frac{\lambda(s,m_A^2,m_{\pi^0}^2)}{4m_A^2}
\qquad\text{(isotropic)},\\
\Big\langle\sum_{\rm pol}|\mathcal M_{\pm\to\pm0}|^2\Big\rangle_\Omega&=\frac{4\Da^2}{F^2}
\Big[\frac{E_k^2E_0^2+p_i^2p_f^2/3}{m_A^2}-m_{\pi^0}^2\Big],
\end{align}
with CM energies $E_k=(s+m_A^2-m_{\pi^\pm}^2)/2\sqrt s$ and $E_0=(s+m_{\pi^0}^2-m_{\pi^\pm}^2)/2\sqrt s$.
Then $\sigma=\tfrac13\langle\sum|\mathcal M|^2\rangle\,p_f/(16\pi s\,p_i)$. Neither final state
has identical particles. The thresholds are
$\sqrt s\ge\max(m_A+m_{\pi^0},2m_{\pi^\pm})$ and $\max(m_A+m_{\pi^\pm},m_{\pi^\pm}+m_{\pi^0})$.

$\Da$ is normalised as for $\Gamma(A'\to3\pi)$ in \texttt{axial3pi\_weight.m}:
$\Da_{\rm eff}=m_{\pi^0A'}/F$ from \texttt{axial\_mixing.m}. The per-$A'$ rate is
\begin{equation}
\Gamma_{A'\pi}(T)=\frac{1}{n_{A'}^{\rm eq}}\sum_{\pi^0,\pi^+,\pi^-}\frac{g_{A'}T}{32\pi^4}\int ds\,
\sigma(s)\,\frac{\lambda(s,m_A^2,m_\pi^2)}{\sqrt s}\,K_1(\sqrt s/T),\qquad
n_{A'}^{\rm eq}=\frac{3m_A^2T}{2\pi^2}K_2(m_A/T).
\end{equation}
This is Maxwell--Boltzmann statistics with $g_{A'}=3$ (\texttt{apipi\_phi\_rate.m}). The pions
are in equilibrium, so detailed balance gives
$dY_{A'}/d\ln x\supset-(\Gamma_{A'\pi}/H)(Y_{A'}-Y_{A',\rm eq})$. It goes into
\texttt{freeze\_coupled}'s $\Gamma_{\rm sc}$ next to $\pi\gamma\to\pi\phi$, so it also enters the
$Y_{A'}$ initial condition. It is on with \texttt{production='sm'}. It is off with
\texttt{LOOP\_PHI\_SC=0} or \texttt{LOOP\_APIPI=0}.

\paragraph{Size.} $\Gamma_{A'\pi}/H\approx4\times10^{11}\eta^2$ at $T=10$ MeV and
$8\times10^{3}\eta^2$ at 5 MeV, for $m_A=0.1$ and $a=\eta[2,-1,-1]$. So the $A'$ reaches the pion
bath at $T_{\rm RH}$ for $\eta\gtrsim2\times10^{-6}$. That is $O(\Da^2)$ against
$O(\alpha^2\Da^4)$ for $A'A'\to\gamma\gamma$, so it is the $A'$'s main bath link wherever
$A'\to3\pi$ is closed. The kaon analogues pay an extra $e^{-(m_K-m_\pi)/T}$ and are dropped.

\section{Numbers ($m_{A'}=0.1$ GeV, $g_\chi=0.01$, $a=\eta[2,-1,-1]$, $T_{\rm RH}=10$ MeV)}

The relic loci are recomputed with the current solver. $\eta$ is where $\Omega_\chi h^2=0.12$
(FIMP root, WIMP root). The first two columns of $\Omega_{A'}h^2$ have no $A'\pi\to\pi\pi$; the
last has every term in.

\begin{center}
\begin{tabular}{cccccc}
\toprule
& \multicolumn{3}{c}{without $A'\pi\to\pi\pi$} & \multicolumn{2}{c}{with $A'\pi\to\pi\pi$}\\
\cmidrule(lr){2-4}\cmidrule(lr){5-6}
$m_\chi$ [GeV] & $\eta$ & $\Omega_{A'}h^2$, no $c_{AA}$ & with $c_{AA}$ & $\eta$ & $\Omega_{A'}h^2$\\
\midrule
0.055 & $2.4\times10^{-3}$ & 0.0038 & 0.0043 & $2.4\times10^{-3}$ & 0.0043\\
0.07 & $1.5\times10^{-6}$, $0.062$ & 0.0014, 0.0056 & 0.0014, 0.0056 & $8.1\times10^{-8}$, 0.062 & 0.0030, 0.0056\\
0.09 & $1.9\times10^{-5}$, $0.24$ & 0.0095, 0.017 & 0.0095, 0.0085 & $7.3\times10^{-8}$, 0.24 & 0.0098, 0.0085\\
0.11 & $3.2\times10^{-4}$, $0.066$ & 1.4, 8.1 & 1.4, 0.99 & $2.4\times10^{-7}$, 0.018 & 1.4, 1.25\\
\bottomrule
\end{tabular}
\end{center}
For $m_\chi<m_{A'}$, the process $A'A'\to\chi\bar\chi$ drains the $A'$ and the remnant is
$\lesssim10\%$ of the DM. For $m_\chi>m_{A'}$, the decay $\chi\bar\chi\to A'A'$ fills it and the
point overcloses even with $A'A'\to$ SM. $c_{AA}$ matters only at large $\eta$, as expected from
$\Da^4$.

$A'\pi\to\pi\pi$ does not change $\Omega_{A'}$ much. It moves the FIMP loci by 1--3 decades in
$\eta$ wherever the $A'$ is stable ($m_\chi>m_{A'}/2$). Below the old FIMP root it already
brings the $A'$ into the pion bath at $T_{\rm RH}$, and $A'A'\to\chi\bar\chi$ ($g_\chi=0.01$) then
makes the $\chi$, independent of $\eta$. The WIMP root at $m_\chi=0.11$ moves from 0.066 to
0.018. At $m_\chi=0.055$, $A'\leftrightarrow\chi\bar\chi$ already ties the $A'$ to the bath
through the $\chi$, and nothing changes.

\section{Not included}

\begin{itemize}
\item $A'K\to K\pi$ and other kaon analogues of $A'\pi\to\pi\pi$, which carry an extra
  $e^{-(m_K-m_\pi)/T}$.
\item The $O(p^4)$ corrections to the $A'\pi\pi\pi$ contact.
\item The anomaly-induced $A'\gamma\gamma^*$ coupling (WZW): Primakoff-like $e\gamma\to eA'$.
\end{itemize}

\end{document}
